\section{引言：从固定输入到序列建模}

在之前的课程中，我们一直讨论的是标准的非循环（non-recurrent）神经网络：固定大小的输入映射到固定大小的输出。这种范式适合图像分类等任务，但真实世界中存在大量\textbf{变长序列}数据——语言、视频、音乐、时间序列等。本讲将引入循环神经网络（Recurrent Neural Networks, RNNs），这是一种专门为序列建模设计的神经网络架构。

\begin{figure}[H]
\centering
\includegraphics[width=0.95\textwidth]{slides-images/slide-004.jpg}
\caption{序列建模的多种范式：从左到右分别为 one-to-one（固定输入输出）、one-to-many（如图像描述）、many-to-one（如视频分类）、many-to-many（输入输出长度可不同，如机器翻译），以及同步的 many-to-many（如逐帧视频分类）\protect\footnotemark}
\end{figure}
\footnotetext{Slides 第5页。}

\begin{importantbox}{序列建模的五种范式}
\begin{itemize}
    \item \textbf{One-to-one}：固定输入 $\rightarrow$ 固定输出（如图像分类）
    \item \textbf{One-to-many}：固定输入 $\rightarrow$ 变长输出（如图像描述生成）
    \item \textbf{Many-to-one}：变长输入 $\rightarrow$ 固定输出（如视频分类、情感分析）
    \item \textbf{Many-to-many}（异步）：变长输入 $\rightarrow$ 变长输出（如机器翻译）
    \item \textbf{Many-to-many}（同步）：每个输入对应一个输出（如逐帧视频标注）
\end{itemize}
\end{importantbox}

\subsection{课程前期回顾}

讲者首先对上节课的内容进行了几点澄清：

\begin{itemize}
    \item \textbf{Dropout 的缩放}：超参数 $p$ 在不同实现中可能表示丢弃概率或保留概率。核心原则是测试时的期望输出应与训练时一致——若训练时丢弃 25\% 的激活值，测试时需乘以 0.75 进行缩放。
    \item \textbf{Layer Normalization 与权重初始化}：Layer Norm 可以缓解不良初始化带来的问题，但无法完全替代良好的权重初始化。在某些需要精确空间信息的任务中，Layer Norm 可能反而损害性能。
\end{itemize}

\subsection{本章小结}

神经网络不再局限于固定大小的输入输出。通过引入序列建模，我们能够处理语言、视频、对话等变长数据。RNN 正是为这类任务设计的核心架构。

%% ========== RNN 基础 ==========

